\section{Related Work}
\label{sec:related}

\subsection{Class-Incremental Learning with Pre-trained Vision-Language Models}

Traditional class-incremental learning (CIL) requires a model to incorporate new classes while retaining discrimination over previously observed ones~\cite{de2021continual,wang2024b}. Representative strategies preserve knowledge through parameter regularization~\cite{kirkpatrick2017overcoming,ahn2019uncertainty}, knowledge distillation~\cite{li2017learning}, rehearsal or experience replay~\cite{rebuffi2017icarl,rolnick2019experience}, and parameter isolation or prompting~\cite{wang2022learning,wang2022dualprompt}. Prototype-based prediction and classifier correction further address the imbalance between old and newly introduced classes. These methods establish historical-task preservation as the central CIL objective, but generally assume that the representation itself has no independent pre-trained capability that must also be retained.

Pre-trained vision-language models extend this objective. In addition to old-class discrimination, continual adaptation should preserve open-vocabulary recognition and the image--text geometry acquired from large-scale pre-training. ZSCL~\cite{zheng2023zscl} constrains sequential adaptation to prevent zero-shot transfer degradation, while Mod-X~\cite{ni2023modx} analyzes intra-modal rotation and inter-modal deviation. C-CLIP~\cite{liu2025cclip} studies multimodal continual learning, and MG-CLIP~\cite{huang2025mgclip} explicitly preserves and compensates for the modality gap. RAIL~\cite{xu2024rail} and LADA~\cite{luo2025lada} further demonstrate the importance of prediction design under task-agnostic cross-domain evaluation. In particular, LADA freezes the visual encoder, adapts text representations, and combines a label-specific visual memory with a text classifier. Our framework instead continually adapts the representation and assigns historical-task preservation and cross-modal preservation to distinct training mechanisms.

\subsection{Knowledge-Preserving Low-Rank Adaptation}

Parameter-efficient fine-tuning adapts pre-trained models through a small set of task-dependent parameters. LoRA~\cite{hu2022lora} constrains residual weight updates to a low-rank factorization, while adapters, prompts~\cite{wang2022sprompts,yu2024moe}, and dynamic-rank variants~\cite{wang2024lora,lu2026codyra} control different aspects of adaptation capacity. Parameter efficiency, however, does not itself specify which historical representation directions may be modified. A low-rank update can still align with high-energy historical activations and substantially change the corresponding responses.

Subspace-based continual-learning methods instead protect historical knowledge by constraining update directions. GPM~\cite{saha2021gradient} projects gradients onto an approximate null space of previous-task features, and DMNSP~\cite{kang2025dynamic} extends multi-layer null-space projection to vision-language continual learning. Null-space-inspired LoRA methods also use low-energy activation directions to initialize or constrain the input-side factor~\cite{tang2026lora}. These approaches act on gradients or on the parameterization of an individual factor. LoRA-NF differs by placing a history-induced filter before both trainable low-rank factors, thereby maintaining the filtered forward form $(\mathbf{X}\mathbf{P})\mathbf{A}\mathbf{B}$ throughout optimization. Cross-modal distillation complements this directional protection by explicitly constraining image--text relations rather than historical activations alone.

\subsection{Continual Prediction with Pre-trained Semantic Priors}

Prototype, nearest-class-mean, and Gaussian discriminant classifiers summarize observed classes without repeatedly optimizing a parametric output layer~\cite{rebuffi2017icarl}. In pre-trained-model CIL, analytic classifiers and label memories can exploit stable feature statistics, whereas the CLIP text classifier provides a language-derived prior over both observed and unobserved classes. RAIL~\cite{xu2024rail} relies on the pre-trained text classifier for transfer, and LADA~\cite{luo2025lada} combines text predictions with a supervised label-memory branch. Gaussian discriminant analysis has also been applied to deep features~\cite{lee2018simple}, but class-specific full covariances are costly and poorly estimated in few-shot regimes. We use LR-RGDA to model regularized class-conditional structure with low-rank covariance corrections, and combine its supervised scores with the CLIP text classifier. This prediction module is evaluated separately from the representation-learning mechanisms.
